% Textos fijos de lang-es.tex que cambian en este juego: va en disco duro o
% tarjeta SD con MSX-DOS2, pide un V9990 y se maneja con rat'on.
\renewcommand{\tCoverLine}{El cl\'asico de rol de PC, en MSX con el V9990}
\renewcommand{\tComputerText}{Un \textbf{MSX2, MSX2+ o MSX turbo R}, con ranuras para el V9990 y la interfaz de disco.}
\renewcommand{\tMemoryText}{\textbf{Al menos \ramNeeded{} de memoria mapeada en total}. Si tu equipo tiene menos, necesita una ampliación, también en los modelos turbo R que no alcancen esa cantidad.}
\renewcommand{\tVideoText}{Un \textbf{cartucho de vídeo Yamaha V9990}. Es obligatorio para la imagen de la aventura.}
\renewcommand{\tSoundText}{Una \textbf{MoonSound} es opcional. Reproduce los efectos y la música de la introducción.}
\renewcommand{\tGameRow}{Disco}
\renewcommand{\tGameText}{\textbf{MSX-DOS2 o Nextor} y un \textbf{disco duro o tarjeta SD/CF} con los archivos del juego.}
\renewcommand{\tControlsText}{Un \textbf{rat\'on MSX en el puerto 1} (se usa para casi todo) y el teclado (moverse y atacar).}
\renewcommand{\tNoTape}{Si te faltan ranuras para cartuchos, usa un \textbf{expansor de ranuras}.}
\renewcommand{\tSetOff}{\textbf{Apaga el MSX.} No insertes ni retires nunca un cartucho con el equipo encendido.}
\renewcommand{\tSetV}{Inserta el \textbf{cartucho V9990} en una ranura libre.}
\renewcommand{\tSetGame}{Inserta la \textbf{interfaz de disco} en otra ranura libre, con los ficheros del juego, y la \textbf{MoonSound} si la tienes. Puedes usar un \textbf{expansor de ranuras} si necesitas conectar más cartuchos.}
\renewcommand{\tSetMonitor}{Conecta el monitor o televisor a la \textbf{salida de v\'ideo del V9990}.}
\renewcommand{\tSetJoy}{Conecta el \textbf{rat\'on al puerto 1}.}
\renewcommand{\tSetOn}{Copia la carpeta \texttt{EOB} de \texttt{SD} dentro del ZIP a tu unidad DOS. Enciende el MSX, selecciona esa unidad y escribe \texttt{CD \textbackslash EOB} y \texttt{EOB}. Elige el idioma en la portada.}
\renewcommand{\tBlackScreen}{\textbf{En openMSX.} Usa un MSX2 o turbo R con la extensión \texttt{gfx9000} y un ratón en el puerto~1. Añade \texttt{moonsound} para oír los efectos y la música de la introducción. El README del proyecto explica cómo arrancar el juego desde el ZIP.\par\medskip
\textbf{Si la pantalla se queda en negro.} Comprueba la conexión del V9990 y selecciona su salida de vídeo. Si aparece el indicador de DOS, comprueba que has arrancado MSX-DOS2 o Nextor y que estás en la carpeta del juego.}
\renewcommand{\tPowerOnCaption}{Izquierda: la portada, donde se elige el idioma. Derecha: la pantalla de t\'itulo con su men\'u.}
\renewcommand{\tJoystick}{Rat\'on}
\renewcommand{\tCreditsTitle}{17. Créditos y derechos}
\renewcommand{\tRightsText}{Basado en las reglas de ADVANCED DUNGEONS \& DRAGONS\textsuperscript{\textregistered} y en el mundo de FORGOTTEN REALMS\textsuperscript{\textregistered}, marcas de sus titulares. El juego, sus gr\'aficos, textos, sonidos y nombre pertenecen a sus titulares. Esta es una versi\'on no oficial hecha para el estudio, la conservaci\'on y como homenaje; no tiene relaci\'on con ellos.}
\newcommand{\tCreditsExtra}{%
\item La portada utiliza la ilustración de apertura elegida para esta edición. Los iconos proceden de los gráficos del juego; varias explicaciones siguen su manual original.
\item Las imágenes del bestiario son capturas de la aventura.
\item ScummVM (licencia GPL) sirvió de referencia para esta adaptación. Su controlador AdLib se usó para convertir el sonido.}
% the pictures of this language: made by tools/manual_images.py (img/) and
% taken by tools/manual_shots.sh (../img/)
\newcommand{\imgdir}{img/}
\newcommand{\shotdir}{../img/}
\newcommand{\tTitleShot}{img/title.png}
\renewcommand{\portWork}{\portWorkES}

\renewcommand{\tMonitorText}{Monitor o TV conectado al V9990 para la aventura. La salida del propio MSX muestra el automapa: usa una segunda pantalla o un monitor con entradas conmutables.}
